\documentclass[../Odyssey_Project.tex]{subfiles}

\begin{document}
\setcounter{section}{4}
\section{Parabolic Subgroups}

\subsubsection*{Setup and Complete Flags}
\begin{itemize}
    \item Throughout, $G = \mathrm{GL}(n, F)$ for some field $F$ and some $n \in \N$, with $V_n(F)$ the space of column vectors of length $n$ over $F$.
    \item Elements of $G$ are viewed as matrices, with respect to the standard basis $\mathbf{v}_1, \dots, \mathbf{v}_n$, of invertible linear transformations of $V_n(F)$.
    \item A \emph{complete flag} on $V_n(F)$ is a sequence of subspaces
    \[
        0 \subset V_1 \subset V_2 \subset \dots \subset V_{n-1} \subset V_n = V_n(F),
    \]
    denoted $(V_1, \dots, V_n)$. Since $\subset$ is proper containment, $\dim_F V_i = i$.
    \item The \emph{standard flag} is given by $V_i = F\mathbf{v}_1 \oplus \dots \oplus F\mathbf{v}_i = V_{i-1} \oplus F\mathbf{v}_i$ (with $V_0 = 0$).
\end{itemize}

\subsubsection*{The $G$-Action on Complete Flags}
\begin{itemize}
    \item $G$ acts on complete flags by $g(V_1, \dots, V_n) = (gV_1, \dots, gV_n)$.
    \item This action is transitive: any complete flag $(W_1, \dots, W_n)$ is obtained from the standard flag by the matrix whose $i$th column is $\mathbf{w}_i \in W_i - W_{i-1}$.
    \item The stabilizer of the standard flag is precisely the standard Borel subgroup $B$, so by Lemma~\ref{lem:3.2} the Borel subgroups of $G$ are exactly the stabilizers of complete flags on $V_n(F)$.
\end{itemize}

\subsubsection*{Upper Unitriangular Matrices and the Decomposition of $B$}
\begin{itemize}
    \item An upper triangular matrix is \emph{upper unitriangular} if all its main-diagonal entries equal $1$, equivalently if it is upper triangular with only eigenvalue $1$.
    \item The set $U$ of all invertible upper unitriangular matrices is a subgroup of $B$ (imitate the proof of Proposition~\ref{prop:4.3}).
    \item Let $T$ be the set of diagonal matrices in $G$. Then $T \leq B$ and $U \cap T = 1$.
\end{itemize}

\begin{proposition}
\label{prop:5.1}
$B = U \rtimes T$.
\end{proposition}
\begin{proof}
Let $\varphi \colon B \to T$ be the map which sends an upper triangular matrix to the diagonal matrix with the same main diagonal. If $b=(b_{ij})$ and $c=(c_{ij})$ lie in $B$, then the $(i,i)$-entry of $bc$ is $b_{ii}c_{ii}$, since all entries below the main diagonal are zero. Hence
\[
    \varphi(bc)=\varphi(b)\varphi(c),
\]
so $\varphi$ is a homomorphism. Its kernel consists exactly of those upper triangular matrices whose diagonal entries are all $1$, so $\ker \varphi=U$. Therefore $U \trianglelefteq B$.\\
Now let $b \in B$. Since $\varphi(b) \in T$ and $\varphi$ restricts to the identity map on $T$, we have
\[
    \varphi(b\varphi(b)^{-1})=\varphi(b)\varphi(b)^{-1}=I.
\]
Thus $b\varphi(b)^{-1} \in U$, and so
\[
    b=(b\varphi(b)^{-1})\varphi(b) \in UT.
\]
Therefore $B=UT$. Finally, $U \cap T=1$, because a diagonal matrix whose diagonal entries are all $1$ is the identity matrix. Hence $B=U \rtimes T$.
\end{proof}

\subsubsection*{Action of $U$ and $T$ on the Standard Flag}
\begin{itemize}
    \item An element of $U$ sends each $\mathbf{v}_i$ to $\mathbf{v}_i$ plus some element of $V_{i-1}$; conversely, any matrix with this property lies in $U$.
    \item Hence $U$ consists of the matrices stabilizing each $V_i$ and inducing the identity on each quotient $V_i / V_{i-1}$.
    \item $T$ consists exactly of the matrices stabilizing each of $F\mathbf{v}_1, \dots, F\mathbf{v}_n$; we say $T$ is the \emph{common stabilizer} of the $F\mathbf{v}_i$.
\end{itemize}

\subsubsection*{Unipotent Elements and Subgroups}
\begin{itemize}
    \item An element of $G$ is \emph{unipotent} if its characteristic polynomial is $(X - 1)^n$.
    \item By Jordan form, any unipotent element of $G$ is conjugate to an element of $U$.
    \item A subgroup of $G$ is \emph{unipotent} if all its elements are unipotent; $U$ is unipotent and comprises all unipotent elements of $B$.
\end{itemize}

\begin{namedtheorem}[Kolchin's Theorem]
\label{thm:kolchin}
Any unipotent subgroup of $G$ is conjugate with a subgroup of $U$.
\end{namedtheorem}
\begin{proof}
Let $H$ be a unipotent subgroup of $G$. We first observe that it is enough to show that $H$ stabilizes some complete flag on $V_n(F)$. Indeed, if $H$ stabilizes a complete flag $\mathcal{F}$, then $\mathcal{F}=g\mathcal{F}_0$ for some $g \in G$, where $\mathcal{F}_0$ is the standard flag. By Lemma~\ref{lem:3.2}, the stabilizer of $\mathcal{F}$ is $gBg^{-1}$, and hence $g^{-1}Hg \leq B$. Since conjugation preserves characteristic polynomials, every element of $g^{-1}Hg$ is still unipotent. The unipotent elements of $B$ all lie in $U$, so $g^{-1}Hg \leq U$. Thus $H$ is conjugate with a subgroup of $U$.\\
It remains to show that $H$ stabilizes some complete flag. We argue by induction on $n$. The case $n=1$ is immediate. Assume $n>1$. By Proposition~\ref{prop:13.28}, there exists some $0 \neq \mathbf{v} \in V_n(F)$ such that $x\mathbf{v}=\mathbf{v}$ for every $x \in H$. Let $W=F\mathbf{v}$. Then $H$ stabilizes $W$, and hence acts on the quotient $V_n(F)/W$. For each $x \in H$, the induced transformation on $V_n(F)/W$ is again unipotent: with respect to a basis beginning with $\mathbf{v}$, the matrix of $x$ has the block form
\[
    \begin{pmatrix}
        1 & * \\
        0 & \overline{x}
    \end{pmatrix},
\]
so the characteristic polynomial of $\overline{x}$ is $(X-1)^{n-1}$. Therefore the induced subgroup on $V_n(F)/W$ is unipotent. By the induction hypothesis, it stabilizes a complete flag on $V_n(F)/W$. Pulling this flag back to $V_n(F)$ and placing $W$ at the beginning gives a complete flag on $V_n(F)$ stabilized by $H$. This completes the induction, and hence proves the theorem.
\end{proof}

\subsubsection*{General Flags and Parabolic Subgroups}
\begin{itemize}
    \item A \emph{flag} on $V_n(F)$ is a nested sequence of non-zero subspaces of $V_n(F)$ which terminates in $V_n(F)$ and has no repeated terms; that is, a sequence $(W_1, \dots, W_r)$ with
    \[
        0 \subset W_1 \subset W_2 \subset \dots \subset W_{r-1} \subset W_r = V_n(F).
    \]
    Since $\subset$ is proper, $r \leq n$; a complete flag is the case $r = n$.
    \item The $G$-set of all flags is not transitive: two flags $(W_1, \dots, W_r)$ and $(W_1', \dots, W_s')$ lie in the same orbit iff $r = s$ and $\dim_F W_i = \dim_F W_i'$ for all $i$ (same \emph{dimension sequence}).
    \item A \emph{parabolic subgroup} of $G$ is the stabilizer of some flag on $V_n(F)$.
    \item Given a flag $(W_1, \dots, W_r)$ with parabolic stabilizer $P$, choose subspaces $Y_i$ with $W_i = W_{i-1} \oplus Y_i$. The \emph{unipotent radical} of $P$ is the subgroup $U_P \leq P$ of matrices inducing the identity on each $W_i / W_{i-1}$; e.g.\ $U_B = U$.
    \item A \emph{Levi complement} of $U_P$ is the subgroup $L_P \leq P$ which is the common stabilizer of the $Y_i$.
    \item A Levi complement is isomorphic with $\mathrm{GL}(y_1, F) \times \dots \times \mathrm{GL}(y_r, F)$, where $y_i = \dim_F Y_i = \dim_F W_i - \dim_F W_{i-1}$. (In particular, any two Levi complements of $U_P$ are isomorphic.)
    \item Taking $Y_i = F\mathbf{v}_i$ gives $L_B = T$, isomorphic with a direct product of $n$ copies of $\mathrm{GL}(1, F) \cong F^\times$.
\end{itemize}

\begin{proposition}
\label{prop:5.2}
If $P$ is a parabolic subgroup of $G$, then we have $P = U_P \rtimes L_P$, where $U_P$ is the unipotent radical of $P$ and $L_P$ is a Levi complement of $U_P$. Furthermore, $P = N_G(U_P)$.
\end{proposition}
\begin{proof}
Let $P$ be the stabilizer of the flag $(W_1,\dots,W_r)$, and put $W_0=0$. Choose subspaces $Y_i$ such that $W_i=W_{i-1}\oplus Y_i$ for each $i$. Then
\[
    V_n(F)=Y_1\oplus \dots \oplus Y_r.
\]
With respect to this direct-sum decomposition, elements of $P$ are block upper triangular. This means that if $p\in P$ and $y_i\in Y_i$, then $p(y_i)$ may have a part in the current block $Y_i$ and earlier blocks $Y_1,\dots,Y_{i-1}$, but it cannot have a part in any later block. Equivalently, there are maps $a_i \colon Y_i\to Y_i$ such that
\[
    p(y_i)=w_{i-1}+a_i(y_i)
\]
with $w_{i-1}\in W_{i-1}$. The map $a_i$ is the action induced by $p$ on the quotient $W_i/W_{i-1}$, so it is invertible.\\
We now split $p$ into its diagonal-block part and its upper-block-unipotent part. Define $l\in G$ by letting $l$ act as $a_i$ on each $Y_i$ and by making $l$ preserve every $Y_i$. Then $l\in L_P$, because $l$ is exactly the block diagonal part of $p$. Set
\[
    u=pl^{-1}.
\]
We claim that $u\in U_P$. Let $y_i\in Y_i$, and write $z_i=l^{-1}(y_i)\in Y_i$. Since $l$ acts as $a_i$ on $Y_i$, we have $a_i(z_i)=y_i$. Hence
\[
    u(y_i)=p(z_i)=y_i+w_{i-1}
\]
for some $w_{i-1}\in W_{i-1}$. Therefore, modulo $W_{i-1}$, the vector $u(y_i)$ is the same as $y_i$. Since the $Y_i$ build the flag one block at a time, this says exactly that $u$ stabilizes every $W_i$ and induces the identity transformation on each quotient $W_i/W_{i-1}$. Thus $u\in U_P$, and
\[
    p=ul\in U_PL_P.
\]
Since $p\in P$ was arbitrary, $P=U_PL_P$.\\
This product is a semidirect product. First, $U_P$ is normal in $P$, since it is the kernel of the homomorphism
\[
    P\to \GL(W_1/W_0)\times \dots \times \GL(W_r/W_{r-1})
\]
which records the induced actions on the successive quotients. Second, $U_P\cap L_P=1$. Indeed, if $x\in U_P\cap L_P$, then $x$ preserves each $Y_i$, but it also acts as the identity on $W_i/W_{i-1}$. Therefore $x$ is the identity on every block $Y_i$, and hence $x=1$. We have proved
\[
    P=U_P\rtimes L_P.
\]
It remains to prove that $P=N_G(U_P)$. Since $U_P\trianglelefteq P$, every element of $P$ normalizes $U_P$, so
\[
    P\leq N_G(U_P).
\]
Conversely, suppose that $g\in N_G(U_P)$. We shall prove that $g$ stabilizes every space in the flag, and hence that $g\in P$.\\
First we recover $W_1$ from $U_P$ alone. We claim that $W_1$ is exactly the common fixed subspace of $U_P$, meaning
\[
    W_1=\{v\in V_n(F)\mid uv=v \text{ for every }u\in U_P\}.
\]
Every vector in $W_1$ is fixed because elements of $U_P$ act as the identity on $W_1/W_0=W_1$. Conversely, take a vector
\[
    v=y_1+\dots+y_r
\]
with $y_i\in Y_i$. If some $y_j\neq0$ for $j>1$, choose a non-zero vector $x\in Y_1$. Extend $y_j$ to a basis of $Y_j$, and define a block-unitriangular element $u$ by
\[
    u(y_j)=y_j+x,
\]
while fixing the remaining basis vectors of $Y_j$ and fixing all other blocks. This element lies in $U_P$: it only adds an earlier-block vector to a vector in the $j$th block, so it acts as the identity on every quotient $W_i/W_{i-1}$. But then
\[
    uv=v+x\neq v.
\]
Thus any vector with a non-zero component outside $W_1$ is moved by some element of $U_P$. Hence the common fixed subspace of $U_P$ is exactly $W_1$.\\
Now use the assumption that $g$ normalizes $U_P$. We show that $gW_1=W_1$. Let $w\in W_1$. To prove $gw\in W_1$, it is enough to show that $gw$ is fixed by every element of $U_P$. Let $u\in U_P$. Since $g\in N_G(U_P)$, we have $g^{-1}ug\in U_P$. Therefore
\[
    u(gw)=g(g^{-1}ug)w.
\]
The element $g^{-1}ug$ lies in $U_P$, and $w\in W_1$ is fixed by every element of $U_P$. Hence $(g^{-1}ug)w=w$, so
\[
    u(gw)=gw.
\]
This holds for every $u\in U_P$, so $gw$ belongs to the common fixed subspace of $U_P$, which is $W_1$. Thus $gW_1\leq W_1$, and equality follows because $g$ is invertible. Therefore $gW_1=W_1$.\\
Now pass to the quotient $V_n(F)/W_1$. On this quotient, the group $U_P$ induces the unipotent radical corresponding to the shorter flag
\[
    W_2/W_1\subset \dots \subset W_r/W_1.
\]
Because $g$ normalizes $U_P$ and stabilizes $W_1$, the induced map of $g$ on $V_n(F)/W_1$ normalizes this induced unipotent radical. The same common-fixed-subspace argument on the quotient shows that the first fixed layer of the quotient, namely $W_2/W_1$, is stabilized by $g$. Hence $gW_2=W_2$. Repeating this argument one layer at a time gives
\[
    gW_i=W_i
\]
for every $i$. Thus $g$ stabilizes the whole flag, so $g\in P$. Therefore $N_G(U_P)\leq P$. Combining both inclusions gives $P=N_G(U_P)$.
\end{proof}

\subsubsection*{Subflags and Staircase Groups}
\begin{itemize}
    \item A \emph{subflag} of the standard flag $(V_1, \dots, V_n)$ is a flag $(W_1, \dots, W_r)$ in which each $W_i = V_j$ for some $j$; there are $2^{n-1}$ such subflags.
    \item A \emph{staircase group} is a parabolic subgroup of $G$ that stabilizes some subflag of the standard flag.
    \item Example ($n = 6$, subflag $0 \subset V_2 \subset V_3 \subset V_6$): the staircase group consists of all matrices in $G = \mathrm{GL}(6, F)$ of the form
    \[
        \begin{pmatrix}
            * & * & * & * & * & * \\
            * & * & * & * & * & * \\
            0 & 0 & * & * & * & * \\
            0 & 0 & 0 & * & * & * \\
            0 & 0 & 0 & * & * & * \\
            0 & 0 & 0 & * & * & *
        \end{pmatrix}.
    \]
    Imagine a ``staircase'' separating the zero entries from the arbitrary entries.
    \item Its unipotent radical consists of all matrices in $G$ of the form
    \[
        \begin{pmatrix}
            1 & 0 & * & * & * & * \\
            0 & 1 & * & * & * & * \\
            0 & 0 & 1 & * & * & * \\
            0 & 0 & 0 & 1 & 0 & 0 \\
            0 & 0 & 0 & 0 & 1 & 0 \\
            0 & 0 & 0 & 0 & 0 & 1
        \end{pmatrix},
    \]
    and the Levi complement (with canonical choices $Y_2 = F\mathbf{v}_3$ and $Y_3 = F\mathbf{v}_4 \oplus F\mathbf{v}_5 \oplus F\mathbf{v}_6$) consists of all matrices in $G$ of the form
    \[
        \begin{pmatrix}
            * & * & 0 & 0 & 0 & 0 \\
            * & * & 0 & 0 & 0 & 0 \\
            0 & 0 & * & 0 & 0 & 0 \\
            0 & 0 & 0 & * & * & * \\
            0 & 0 & 0 & * & * & * \\
            0 & 0 & 0 & * & * & *
        \end{pmatrix},
    \]
    isomorphic with $\mathrm{GL}(2, F) \times \mathrm{GL}(1, F) \times \mathrm{GL}(3, F)$.
    \item In general, the unipotent radical of any staircase group is contained in $B$.
    \item The $G$-set of all flags is partitioned into the orbits of the subflags of the standard flag, since any flag has the same dimension sequence as some subflag of the standard flag. By Lemma~\ref{lem:3.2}, the parabolic subgroups of $G$ are exactly the conjugates of the staircase groups.
\end{itemize}

\subsubsection*{Subgroups Containing $B$}


\begin{lemma}
\label{lem:5.3}
The only non-zero subspaces of $V_n(F)$ left invariant by $B$ are the subspaces $V_1, \dots, V_n$ that appear in the standard flag.
\end{lemma}
\begin{proof}
First observe that each $V_i$ is stabilized by $B$: if $j\leq i$ and $b\in B$, then $b\mathbf{v}_j$ is a linear combination of $\mathbf{v}_1,\dots,\mathbf{v}_j$, hence lies in $V_i$. Therefore $bV_i\subseteq V_i$, and equality follows since $b$ is invertible.\\
Conversely, let $Y$ be a non-zero subspace of $V_n(F)$ stabilized by $B$. Choose $k$ minimal such that $Y\subseteq V_k$. Then $Y\not\subseteq V_{k-1}$, so there is some
\[
    y=\alpha_1\mathbf{v}_1+\dots+\alpha_k\mathbf{v}_k\in Y
\]
with $\alpha_k\neq 0$. Let $b\in B$ be the matrix which is the identity except that its $k$th column is $y$:
\[
    b=
    \begin{pmatrix}
        I_{k-1} & \begin{matrix}\alpha_1\\ \vdots\\ \alpha_{k-1}\end{matrix} & 0\\
        0 & \alpha_k & 0\\
        0 & 0 & I_{n-k}
    \end{pmatrix}.
\]
This matrix is upper triangular and invertible, since its $k$th diagonal entry is $\alpha_k\neq 0$. Also $b\mathbf{v}_k=y$. Since $Y$ is stabilized by $B$ and $b^{-1}\in B$, we have
\[
    \mathbf{v}_k=b^{-1}y\in Y.
\]
Now let
\[
    z=\beta_1\mathbf{v}_1+\dots+\beta_k\mathbf{v}_k\in V_k-V_{k-1}.
\]
Then $\beta_k\neq 0$, and the same matrix construction gives some $c\in B$ with $c\mathbf{v}_k=z$. Since $\mathbf{v}_k\in Y$ and $Y$ is stabilized by $B$, it follows that $z\in Y$. Thus every vector in $V_k-V_{k-1}$ lies in $Y$. In particular, if $w\in V_{k-1}$, then $w+\mathbf{v}_k\in Y$, and since $\mathbf{v}_k\in Y$, we get $w\in Y$. Hence $V_{k-1}\subseteq Y$, and together with $\mathbf{v}_k\in Y$ this gives $V_k\subseteq Y$. Since $Y\subseteq V_k$ by choice of $k$, we conclude that $Y=V_k$.
\end{proof}

\begin{theorem}
\label{thm:5.4}
The only subgroups of $G$ that contain $B$ are the staircase groups.
\end{theorem}
\begin{proof}
First observe that every staircase group contains $B$, because $B$ stabilizes every subspace in the standard flag and therefore stabilizes every subflag of the standard flag. We prove the converse. Let $H$ be a subgroup with
\[
    B \leq H \leq G.
\]
If a subspace of $V_n(F)$ is stabilized by $H$, then it is also stabilized by $B$. Hence Lemma~\ref{lem:5.3} shows that the only possible non-zero $H$-invariant subspaces are among
\[
    V_1,\dots,V_n.
\]
List the ones which are actually stabilized by $H$ as
\[
    V_{a_1},\dots,V_{a_r},
    \qquad
    1\leq a_1<\dots<a_r=n,
\]
and put $a_0=0$. Let $P$ be the staircase group stabilizing the subflag
\[
    0\subset V_{a_1}\subset\dots\subset V_{a_r}=V_n(F).
\]
Then $H\leq P$. We shall prove that $H=P$.\\
First we prove the basic case in which the only $H$-invariant subspaces are $0$ and $V_n(F)$. In this case we claim that $H=G$. If $n=1$ this is immediate, so assume $n>1$. The subspace
\[
    Y=\langle h\mathbf{v}_1\mid h\in H\rangle
\]
is non-zero and $H$-invariant, so by assumption $Y=V_n(F)$. Thus some $h\in H$ sends $\mathbf{v}_1$ to a vector
\[
    h\mathbf{v}_1=\alpha_1\mathbf{v}_1+\dots+\alpha_n\mathbf{v}_n
\]
with $\alpha_n\neq 0$. By \nameref{thm:bruhat}, write $h=b_1wb_2$ with $b_1,b_2\in B$ and $w\in W$. Since $B\leq H$ and $h\in H$, it follows that $w\in H$. By Corollary~\ref{cor:4.10}, the fact that the last non-zero term in $h\mathbf{v}_1$ is the $\mathbf{v}_n$-term implies that $w(1)=n$. Let $j=w^{-1}(1)$. Then $j>1$, so $X_{1j}\leq B\leq H$. Using Lemma~\ref{lem:4.6}(v), we get
\[
    wX_{1j}(\alpha)w^{-1}=X_{n1}(\alpha)
\]
for every $\alpha\in F$. Hence $X_{n1}\leq H$.\\
We now get all the remaining root subgroups from this one lower root subgroup. The upper root subgroups $X_{ij}$ with $i<j$ already lie in $B$, and hence in $H$. If $n=2$, this already gives every root subgroup. If $n>2$, then for $1<i<n$ and $\alpha\in F$, Lemma~\ref{lem:4.6}(iv) gives
\[
    X_{ni}(\alpha)=[X_{n1}(\alpha),X_{1i}(1)]\in H
\]
and
\[
    X_{i1}(\alpha)=[X_{in}(\alpha),X_{n1}(1)]\in H.
\]
For distinct $1<i,j<n$, the same commutator formula gives
\[
    X_{ij}(\alpha)=[X_{i1}(\alpha),X_{1j}(1)]\in H.
\]
Thus $H$ contains every root subgroup $X_{ij}$, and it also contains every invertible diagonal matrix because all diagonal matrices lie in $B$. By Theorem~\ref{thm:4.12}, $H=G$. This proves the basic case.\\
We return to the general subgroup $H$. By Proposition~\ref{prop:5.2}, the staircase group $P$ decomposes as
\[
    P=U_P\rtimes L_P.
\]
Since $P$ is a staircase group, its unipotent radical $U_P$ is contained in $B$, so $U_P\leq H$. Therefore it is enough to show that a Levi complement $L_P$ is contained in $H$. We choose the standard Levi complement whose blocks are
\[
    F\mathbf{v}_{a_{k-1}+1}\oplus\dots\oplus F\mathbf{v}_{a_k},
    \qquad
    1\leq k\leq r.
\]
Fix one block, and write
\[
    s=a_{k-1}+1,
    \qquad
    t=a_k.
\]
If $s=t$, there is nothing to prove for this block. Assume $s<t$. Because there is no $H$-invariant standard flag subspace strictly between $V_{a_{k-1}}$ and $V_{a_k}$, the subspace
\[
    V_{a_{k-1}}+\langle h\mathbf{v}_s\mid h\in H\rangle
\]
must be all of $V_t$. Hence there is some $h\in H$ such that the coefficient of $\mathbf{v}_t$ in $h\mathbf{v}_s$ is non-zero. Write $h=b_1wb_2$ by \nameref{thm:bruhat}. Again $w\in H$. Also $w\in P$, because $b_1,b_2\in B\leq P$ and $h\in P$. Thus $w$ preserves each set of indices
\[
    \{1,\dots,a_\ell\}
\]
which comes from the subflag stabilized by $P$. In particular, $w$ sends the block $\{s,\dots,t\}$ to itself. The row-pivot description in Corollary~\ref{cor:4.10} applied to the $s$th column says that $w(s)=t$, because the first $s-1$ columns use only rows $1,\dots,s-1$, while the $s$th column has no entries below row $t$ and has a non-zero entry in row $t$. Let $q=w^{-1}(s)$. Since $w$ preserves the block and $w(s)=t$, we have $s<q\leq t$. Therefore $X_{sq}\leq B\leq H$, and Lemma~\ref{lem:4.6}(v) gives
\[
    wX_{sq}(\alpha)w^{-1}=X_{ts}(\alpha)
\]
for every $\alpha\in F$. Hence $X_{ts}\leq H$.\\
Now the same commutator ladder used in the basic case works inside this block. All upper root subgroups $X_{ij}$ with $s\leq i<j\leq t$ lie in $B$, hence in $H$. Together with $X_{ts}\leq H$, this gives every root subgroup inside the block: for $s<i<t$,
\[
    X_{ti}(\alpha)=[X_{ts}(\alpha),X_{si}(1)]\in H
\]
and
\[
    X_{is}(\alpha)=[X_{it}(\alpha),X_{ts}(1)]\in H,
\]
while for distinct $s<i,j<t$,
\[
    X_{ij}(\alpha)=[X_{is}(\alpha),X_{sj}(1)]\in H.
\]
Thus, for each block of $L_P$, the subgroup $H$ contains all root subgroups supported inside that block. Since $H$ also contains all diagonal matrices, Theorem~\ref{thm:4.12} applied block-by-block shows that $L_P\leq H$. Hence
\[
    P=U_P\rtimes L_P\leq H.
\]
We already had $H\leq P$, so $H=P$. Therefore every subgroup of $G$ containing $B$ is a staircase group.
\end{proof}

\subsection*{Exercises}

Let $G = \mathrm{GL}(n, F)$, where $n \geq 2$.

\begin{exercise}
\label{ex:5.1}
Suppose that $g \in G$ fixes some $0 \neq \mathbf{v} \in V_n(F)$ and induces the identity transformation on $V_n(F) / F\mathbf{v}$. Show that $g$ is conjugate with an element of the root subgroup $X_{12}$.
\end{exercise}
\begin{solution}
Since $g$ induces the identity transformation on $V_n(F)/F\mathbf{v}$, for every $\mathbf{x}\in V_n(F)$ we have
\[
    g\mathbf{x}+F\mathbf{v}=\mathbf{x}+F\mathbf{v}.
\]
Equivalently,
\[
    g\mathbf{x}-\mathbf{x}\in F\mathbf{v}.
\]
Thus there is a scalar $\lambda(\mathbf{x})\in F$ such that
\[
    g\mathbf{x}=\mathbf{x}+\lambda(\mathbf{x})\mathbf{v}.
\]
The scalar $\lambda(\mathbf{x})$ is unique because $\mathbf{v}\neq 0$. The assignment $\lambda \colon V_n(F)\to F$ is linear: for example, comparing
\[
    g(\mathbf{x}+\mathbf{y})=(\mathbf{x}+\mathbf{y})+\lambda(\mathbf{x}+\mathbf{y})\mathbf{v}
\]
with
\[
    g\mathbf{x}+g\mathbf{y}
    =\mathbf{x}+\lambda(\mathbf{x})\mathbf{v}
     +\mathbf{y}+\lambda(\mathbf{y})\mathbf{v}
\]
gives $\lambda(\mathbf{x}+\mathbf{y})=\lambda(\mathbf{x})+\lambda(\mathbf{y})$, and scalar multiplication is similar. Also $g\mathbf{v}=\mathbf{v}$, so $\lambda(\mathbf{v})=0$.\\
If $\lambda=0$, then $g\mathbf{x}=\mathbf{x}$ for every $\mathbf{x}$, so $g=I=X_{12}(0)$, and there is nothing to prove. Assume now that $\lambda\neq 0$. Choose $\mathbf{w}_1=\mathbf{v}$. Since $\lambda\neq 0$, there exists some vector $\mathbf{w}_2$ with
\[
    \lambda(\mathbf{w}_2)=\alpha\neq 0.
\]
Then
\[
    g\mathbf{w}_2=\mathbf{w}_2+\alpha\mathbf{w}_1.
\]
The kernel of $\lambda$ has dimension $n-1$. Since $\lambda(\mathbf{w}_1)=0$, we have $\mathbf{w}_1\in\ker\lambda$. Choose vectors $\mathbf{w}_3,\dots,\mathbf{w}_n$ such that
\[
    \mathbf{w}_1,\mathbf{w}_3,\dots,\mathbf{w}_n
\]
is a basis of $\ker\lambda$. Since $\mathbf{w}_2\notin\ker\lambda$, the list $\mathbf{w}_1,\mathbf{w}_2,\dots,\mathbf{w}_n$ is a basis of $V_n(F)$. In this basis, $g$ sends
\[
    \mathbf{w}_1\mapsto\mathbf{w}_1,\qquad
    \mathbf{w}_2\mapsto\mathbf{w}_2+\alpha\mathbf{w}_1,\qquad
    \mathbf{w}_i\mapsto\mathbf{w}_i \quad (i\geq 3).
\]
Therefore the matrix of $g$ in the basis $\mathbf{w}_1,\dots,\mathbf{w}_n$ is $X_{12}(\alpha)$. If $p\in G$ is the change-of-basis matrix whose columns are $\mathbf{w}_1,\dots,\mathbf{w}_n$, then
\[
    p^{-1}gp=X_{12}(\alpha).
\]
Hence $g$ is conjugate with an element of the root subgroup $X_{12}$.
\end{solution}

\begin{exercise}
\label{ex:5.2}
We say a basis $\mathcal{B}$ of $V_n(F)$ \emph{belongs} to a complete flag $(V_1, \dots, V_n)$ if each $V_i$ contains exactly one element of $\mathcal{B}$ that is not in $V_{i-1}$.
\begin{enumerate}
    \item[(a)] Show that if $(V_1, \dots, V_n)$ and $(W_1, \dots, W_n)$ are complete flags on $V_n(F)$, then there is a basis of $V_n(F)$ that belongs to both flags.
    \item[(b)] Use part (a) to give a different proof of the \nameref{thm:bruhat}.
\end{enumerate}
\end{exercise}
\begin{solution}
    \phantom{}\par
    \begin{enumerate}
        \item[(a)] Put $V_0=W_0=0$. We build the required basis one vector at a time from intersections of the two flags. Fix $i\in\{1,\dots,n\}$. Since $W_n=V_n(F)$, we have
    \[
        V_i\cap W_n=V_i,
    \]
    and this is not contained in $V_{i-1}$. Therefore there is a smallest integer $j_i$ such that
    \[
        V_i\cap W_{j_i}\not\leq V_{i-1}.
    \]
    Choose
    \[
        \mathbf{b}_i\in (V_i\cap W_{j_i})-V_{i-1}.
    \]
    This means that $\mathbf{b}_i$ enters the $V$-flag at the $i$th step, and it enters the $W$-flag no later than the $j_i$th step.\\
    First we check that these vectors form a basis belonging to the $V$-flag. Since $\mathbf{b}_1\in V_1$ and $\mathbf{b}_1\neq0$, the first vector works. Suppose that $\mathbf{b}_1,\dots,\mathbf{b}_{i-1}$ are independent and span $V_{i-1}$. Since $\mathbf{b}_i\notin V_{i-1}$, the vector $\mathbf{b}_i$ is not in the span of the previous vectors. Hence $\mathbf{b}_1,\dots,\mathbf{b}_i$ are independent. They are $i$ independent vectors inside the $i$-dimensional space $V_i$, so they span $V_i$. By induction,
    \[
        V_i=\langle \mathbf{b}_1,\dots,\mathbf{b}_i\rangle
    \]
    for every $i$. Thus the basis belongs to the $V$-flag.\\
    It remains to show that the same basis also belongs to the $W$-flag. Fix $m\in\{1,\dots,n\}$. For each $i$, the space
    \[
        V_{i-1}\cap W_m
    \]
    is contained in $V_i\cap W_m$, and the dimension can increase by at most $1$ because $V_i/V_{i-1}$ is one-dimensional. By the definition of $j_i$, this increase happens exactly when $j_i\leq m$. Therefore
    \[
        \#\{i\mid j_i\leq m\}
        =
        \sum_{i=1}^n
        \bigl(\dim(V_i\cap W_m)-\dim(V_{i-1}\cap W_m)\bigr).
    \]
    The sum telescopes:
    \[
        \sum_{i=1}^n
        \bigl(\dim(V_i\cap W_m)-\dim(V_{i-1}\cap W_m)\bigr)
        =
        \dim(V_n\cap W_m)-\dim(V_0\cap W_m)
        =
        \dim W_m
        =
        m.
    \]
    Also, by the minimality of $j_i$, the vector $\mathbf{b}_i$ lies in $W_m$ exactly when $j_i\leq m$. Thus exactly $m$ of the vectors $\mathbf{b}_1,\dots,\mathbf{b}_n$ lie in $W_m$. These $m$ vectors are independent and lie inside the $m$-dimensional space $W_m$, so they form a basis of $W_m$. Hence, when we pass from $W_{m-1}$ to $W_m$, exactly one new basis vector appears. Therefore the same basis belongs to the $W$-flag as well.

        \item[(b)] Let
    \[
        \mathcal{F}_0=(V_1,\dots,V_n)
    \]
    be the standard complete flag, and let $g\in G$. Then
    \[
        g\mathcal{F}_0=(gV_1,\dots,gV_n)
    \]
    is also a complete flag. By part (a), there is a basis
    \[
        \mathcal{B}=\{\mathbf{b}_1,\dots,\mathbf{b}_n\}
    \]
    which belongs to both $\mathcal{F}_0$ and $g\mathcal{F}_0$. Order the basis so that
    \[
        \mathbf{b}_i\in V_i-V_{i-1}
    \]
    for every $i$. Let $a$ be the matrix whose $i$th column is $\mathbf{b}_i$. Since the $i$th column of $a$ lies in $V_i$, the matrix $a$ is upper triangular. Since its columns form a basis, $a$ is invertible. Hence
    \[
        a\in B.
    \]
    Now use the fact that the same basis also belongs to $g\mathcal{F}_0$. For each $r$, exactly one basis vector lies in
    \[
        gV_r-gV_{r-1}.
    \]
    Thus there is a permutation $\sigma$ of $\{1,\dots,n\}$ such that
    \[
        \mathbf{b}_{\sigma(r)}\in gV_r-gV_{r-1}
    \]
    for every $r$. Let $w\in W$ be the permutation matrix which reorders the columns of $a$ so that the $r$th column of $aw$ is $\mathbf{b}_{\sigma(r)}$. Then
    \[
        g^{-1}\mathbf{b}_{\sigma(r)}\in V_r-V_{r-1}
    \]
    for every $r$. Therefore the $r$th column of $g^{-1}aw$ lies in $V_r$. It follows that $g^{-1}aw$ is upper triangular. It is also invertible, so
    \[
        g^{-1}aw\in B.
    \]
    Write
    \[
        c=g^{-1}aw\in B.
    \]
    Then
    \[
        aw=gc,
    \]
    and therefore
    \[
        g=awc^{-1}.
    \]
    Since $a\in B$, $w\in W$, and $c^{-1}\in B$, we have $g\in BWB$. As $g\in G$ was arbitrary, this gives
    \[
        G\subseteq BWB.
    \]
    The reverse inclusion $BWB\subseteq G$ is immediate from $B\leq G$ and $W\leq G$. Hence
    \[
        G=BWB,
    \]
    which is the \nameref{thm:bruhat}.
    \end{enumerate}
\end{solution}

\begin{exercise}
\label{ex:5.3}
Let $P$ and $Q$ be distinct parabolic subgroups of $G$ containing the standard Borel subgroup $B$ of $G$ (so that, by Theorem~\ref{thm:5.4}, $P$ and $Q$ are staircase groups).
\begin{enumerate}
    \item[(a)] Show that $P$ and $Q$ are not conjugate in $G$.
    \item[(b)] If $L_P$ and $L_Q$ are corresponding Levi complements, find conditions on $P$ and $Q$ that determine when $L_P$ and $L_Q$ will be conjugate in $G$.
\end{enumerate}
\end{exercise}
\begin{solution}
    \phantom{}\par
    \begin{enumerate}
        \item[(a)] By Theorem~\ref{thm:5.4}, both $P$ and $Q$ are staircase groups. Thus there are subflags of the standard flag
    \[
        0\subset V_{a_1}\subset \dots \subset V_{a_r}=V_n(F)
    \]
    and
    \[
        0\subset V_{b_1}\subset \dots \subset V_{b_s}=V_n(F)
    \]
    such that $P$ is the stabilizer of the first subflag and $Q$ is the stabilizer of the second. If the two lists of indices were the same, then the two stabilized standard subflags would be the same, and hence their stabilizers would be the same. Since $P\neq Q$, we know that
    \[
        \{a_1,\dots,a_r\}\neq \{b_1,\dots,b_s\}.
    \]
    We first record the invariant-subspace fact we need. Let $R$ be the staircase group stabilizing
    \[
        0\subset V_{c_1}\subset \dots \subset V_{c_t}=V_n(F),
    \]
    and put $c_0=0$. We claim that the only non-zero subspaces of $V_n(F)$ stabilized by $R$ are
    \[
        V_{c_1},\dots,V_{c_t}.
    \]
    This is the block version of Lemma~\ref{lem:5.3}. Let $Y$ be a non-zero $R$-invariant subspace. Choose $k$ minimal such that
    \[
        Y\leq V_{c_k}.
    \]
    Then $Y\not\leq V_{c_{k-1}}$, so the image of $Y$ in the quotient
    \[
        V_{c_k}/V_{c_{k-1}}
    \]
    is non-zero. By Proposition~\ref{prop:5.2} and the block description of a staircase group, the standard Levi complement of $R$ acts as the full general linear group on the block
    \[
        F\mathbf{v}_{c_{k-1}+1}\oplus \dots \oplus F\mathbf{v}_{c_k}.
    \]
    Hence it acts on $V_{c_k}/V_{c_{k-1}}$ as the full group $\GL(c_k-c_{k-1},F)$. Since $Y$ is $R$-invariant, the image of $Y$ in this quotient is invariant under this full general linear group. But the only non-zero subspace invariant under the full general linear group is the whole quotient, because an invertible linear map can send any non-zero vector to any other non-zero vector. Therefore
    \[
        Y+V_{c_{k-1}}=V_{c_k}.
    \]
    We next show that $V_{c_{k-1}}\leq Y$. If $k=1$, this is immediate because $V_{c_0}=0$. Assume $k>1$. Since $Y+V_{c_{k-1}}=V_{c_k}$, there is some vector $y\in Y$ whose component in the $k$th block is non-zero. Choose an index $m$ in the $k$th block such that the coefficient of $\mathbf{v}_m$ in $y$ is a non-zero scalar $\beta$. Now let $x\in V_{c_{k-1}}$. Write
    \[
        x=\gamma_1\mathbf{v}_1+\dots+\gamma_{c_{k-1}}\mathbf{v}_{c_{k-1}}.
    \]
    For each $\ell\leq c_{k-1}$, the root subgroup element $X_{\ell m}(\alpha)$ sends $\mathbf{v}_m$ to $\mathbf{v}_m+\alpha\mathbf{v}_\ell$. Therefore, if we choose
    \[
        u=\prod_{\ell=1}^{c_{k-1}}X_{\ell m}(\gamma_\ell/\beta),
    \]
    then $u$ sends $y$ to $y+x$. These are upper root subgroup elements because $\ell<m$, so they lie in $B$. Since $B\leq R$, we have $u\in R$. Since $Y$ is $R$-invariant, both $y$ and $uy=y+x$ lie in $Y$, and hence
    \[
        x=(y+x)-y\in Y.
    \]
    Since $x\in V_{c_{k-1}}$ was arbitrary, $V_{c_{k-1}}\leq Y$. Together with $Y+V_{c_{k-1}}=V_{c_k}$, this gives $V_{c_k}\leq Y$. Since $Y\leq V_{c_k}$ by the choice of $k$, we get $Y=V_{c_k}$. This proves the claim.\\
    Now suppose, for a contradiction, that $P$ and $Q$ are conjugate in $G$. Then there exists some $g\in G$ such that
    \[
        Q=gPg^{-1}.
    \]
    For each $i$, the subspace $V_{a_i}$ is stabilized by $P$. We claim that $gV_{a_i}$ is stabilized by $Q$. Indeed, if $q\in Q$, then $q=gpg^{-1}$ for some $p\in P$. If $gv\in gV_{a_i}$ with $v\in V_{a_i}$, then
    \[
        q(gv)=gpg^{-1}gv=gpv.
    \]
    Since $p$ stabilizes $V_{a_i}$, we have $pv\in V_{a_i}$, and hence $gpv\in gV_{a_i}$. Thus $gV_{a_i}$ is $Q$-invariant. By the claim applied to $Q$, the subspace $gV_{a_i}$ must be one of the subspaces
    \[
        V_{b_1},\dots,V_{b_s}.
    \]
    Taking dimensions gives
    \[
        a_i=\dim gV_{a_i}=b_j
    \]
    for some $j$. Therefore
    \[
        \{a_1,\dots,a_r\}\subseteq \{b_1,\dots,b_s\}.
    \]
    Applying the same argument to $P=g^{-1}Qg$ gives the reverse inclusion. Hence
    \[
        \{a_1,\dots,a_r\}=\{b_1,\dots,b_s\},
    \]
    contradicting the fact that these two sets are different. Therefore no such $g$ exists, and $P$ and $Q$ are not conjugate in $G$.

        \item[(b)] Keep the notation from part (a), and put $a_0=b_0=0$. The standard Levi complement of $P$ is the subgroup which preserves each block
    \[
        Y_i=F\mathbf{v}_{a_{i-1}+1}\oplus \dots \oplus F\mathbf{v}_{a_i},
    \]
    and acts on that block as the full group $\GL(Y_i)$. Thus its block sizes are
    \[
        d_i=\dim Y_i=a_i-a_{i-1}
        \qquad (1\leq i\leq r).
    \]
    Similarly, the standard Levi complement of $Q$ has blocks
    \[
        Z_j=F\mathbf{v}_{b_{j-1}+1}\oplus \dots \oplus F\mathbf{v}_{b_j}
    \]
    of sizes
    \[
        e_j=\dim Z_j=b_j-b_{j-1}
        \qquad (1\leq j\leq s).
    \]
    We claim that $L_P$ and $L_Q$ are conjugate in $G$ exactly when the two lists
    \[
        d_1,\dots,d_r
        \qquad\text{and}\qquad
        e_1,\dots,e_s
    \]
    agree up to permutation.\\
    First suppose that the block sizes agree up to permutation. Then $r=s$, and after reordering the $Q$-blocks there is a bijection $\sigma$ such that
    \[
        d_i=e_{\sigma(i)}
    \]
    for every $i$. Choose a permutation matrix $w$ which sends the standard basis vectors in $Y_i$ onto the standard basis vectors in $Z_{\sigma(i)}$. Conjugating by $w$ sends an arbitrary invertible linear transformation on $Y_i$ to an arbitrary invertible linear transformation on $Z_{\sigma(i)}$. Therefore
    \[
        wL_Pw^{-1}=L_Q.
    \]
    Hence $L_P$ and $L_Q$ are conjugate.\\
    Conversely, suppose that $L_Q=gL_Pg^{-1}$ for some $g\in G$. Conjugation by $g$ sends the block decomposition
    \[
        V_n(F)=Y_1\oplus \dots \oplus Y_r
    \]
    to
    \[
        V_n(F)=gY_1\oplus \dots \oplus gY_r.
    \]
    On each $gY_i$, the conjugate subgroup $gL_Pg^{-1}=L_Q$ contains a factor acting as the full group $\GL(gY_i)$. Thus the spaces $gY_i$ are the same kind of full general-linear blocks for $L_Q$ as the standard blocks $Z_j$. The standard block decomposition of $L_Q$ has block dimensions $e_1,\dots,e_s$, so the dimensions
    \[
        \dim gY_i=\dim Y_i=d_i
    \]
    must be exactly the same dimensions as the $e_j$, possibly in a different order. Hence the block sizes of $P$ and $Q$ agree up to permutation.\\
    Therefore the condition is:
    \[
        L_P\text{ and }L_Q\text{ are conjugate in }G
        \quad\Longleftrightarrow\quad
        \{a_i-a_{i-1}\mid 1\leq i\leq r\}
        =
        \{b_j-b_{j-1}\mid 1\leq j\leq s\}
    \]
    as multisets.
    \end{enumerate}
\end{solution}

\begin{exercise}
\label{ex:5.4}
Prove Proposition~\ref{prop:5.2}.
\end{exercise}
\begin{solution}
This is exactly Proposition~\ref{prop:5.2}, proved above.
\end{solution}

\begin{exercise}
\label{ex:5.5}
Prove the following generalization of Lemma~\ref{lem:5.3}: If $P$ is a parabolic subgroup of $G$ which is the stabilizer of the flag $(W_1, \dots, W_r)$, then the $W_i$ are the only subspaces of $V_n(F)$ left invariant by $P$.
\end{exercise}
\begin{solution}
Let
\[
    0\subset W_1\subset \dots \subset W_r=V_n(F)
\]
be the flag stabilized by $P$, and put $c_i=\dim_F W_i$. By the transitivity of the $G$-action on flags with a fixed dimension sequence, choose $g\in G$ such that
\[
    gW_i=V_{c_i}
\]
for every $i$. Hence $gPg^{-1}$ is the staircase group stabilizing the standard subflag
\[
    0\subset V_{c_1}\subset \dots \subset V_{c_r}=V_n(F).
\]\\
Now let $Y$ be a non-zero subspace of $V_n(F)$ stabilized by $P$. Then $gY$ is stabilized by $gPg^{-1}$: if $q=gpg^{-1}\in gPg^{-1}$ and $gy\in gY$, with $p\in P$ and $y\in Y$, then
\[
    q(gy)=gpy\in gY.
\]
By the staircase-group claim proved in Exercise~\ref{ex:5.3}(a), the only non-zero subspaces stabilized by $gPg^{-1}$ are
\[
    V_{c_1},\dots,V_{c_r}.
\]
Therefore $gY=V_{c_i}$ for some $i$. Applying $g^{-1}$ gives
\[
    Y=g^{-1}V_{c_i}=W_i.
\]
Thus the only non-zero subspaces stabilized by $P$ are $W_1,\dots,W_r$. Each $W_i$ is stabilized by $P$ by the definition of $P$ as the stabilizer of the flag, so these are exactly the non-zero $P$-invariant subspaces.
\end{solution}

\begin{exercise}
\label{ex:5.6}
(cont.) Using Exercise~\ref{ex:5.5}, show that any parabolic subgroup of $G$ is self-normalizing.
\end{exercise}
\begin{solution}
Let $P$ be the stabilizer of the flag
\[
    0\subset W_1\subset \dots\subset W_r=V_n(F).
\]
We prove that $N_G(P)=P$. The inclusion
\[
    P\leq N_G(P)
\]
is automatic, because every element of a subgroup normalizes that subgroup. It remains to prove the reverse inclusion.\\
Let $g\in N_G(P)$. We want to show that $g\in P$, i.e. that $g$ stabilizes every $W_i$. Fix $i$. We first show that $gW_i$ is stabilized by $P$. Let $p\in P$ and let $gw\in gW_i$, where $w\in W_i$. Since $g\in N_G(P)$, we have
\[
    g^{-1}pg\in P.
\]
Since $P$ stabilizes $W_i$, this gives
\[
    (g^{-1}pg)w\in W_i.
\]
Applying $g$ to both sides, we get
\[
    p(gw)=g(g^{-1}pg)w\in gW_i.
\]
Thus $gW_i$ is a non-zero subspace stabilized by $P$. By Exercise~\ref{ex:5.5}, the only non-zero subspaces stabilized by $P$ are exactly the spaces in its flag, so
\[
    gW_i=W_j
\]
for some $j$. But $g$ preserves dimension, so
\[
    \dim_F W_j=\dim_F(gW_i)=\dim_F W_i.
\]
The dimensions in a flag are strictly increasing, so $j=i$. Hence
\[
    gW_i=W_i
\]
for every $i$. Therefore $g$ stabilizes the whole flag, and so $g\in P$. This proves
\[
    N_G(P)\leq P.
\]
Combining the two inclusions gives $N_G(P)=P$, so every parabolic subgroup of $G$ is self-normalizing.
\end{solution}

\begin{exercise}
\label{ex:5.7}
Complete the outlined proof of Theorem~\ref{thm:5.4}.
\end{exercise}
\begin{solution}
Let $H$ be a subgroup of $G$ with $B\leq H\leq G$. If a subspace is stabilized by $H$, then it is also stabilized by $B$. Hence Lemma~\ref{lem:5.3} says that the only possible non-zero subspaces stabilized by $H$ are among the standard flag subspaces. List the ones actually stabilized by $H$ as
\[
    V_{a_1},\dots,V_{a_r},
    \qquad
    1\leq a_1<\dots<a_r=n,
\]
and put $a_0=0$. Let $P$ be the staircase group stabilizing the subflag
\[
    0\subset V_{a_1}\subset \dots\subset V_{a_r}=V_n(F).
\]
Then $H\leq P$, because every element of $H$ stabilizes each of the spaces in this subflag. We prove the reverse inclusion.\\
\textit{Step 1: reduce to the Levi complement.} By Proposition~\ref{prop:5.2}, we may write
\[
    P=U_P\rtimes L_P.
\]
Since $P$ is a staircase group, its unipotent radical $U_P$ is contained in $B$. Hence
\[
    U_P\leq B\leq H.
\]
Therefore, to prove $P\leq H$, it remains only to prove
\[
    L_P\leq H.
\]
Choose the standard Levi complement whose blocks are
\[
    F\mathbf{v}_{a_{k-1}+1}\oplus\dots\oplus F\mathbf{v}_{a_k},
    \qquad
    1\leq k\leq r.
\]
We will show that $H$ contains the full general linear group on each one of these blocks.\\
\textit{Step 2: find one lower root subgroup in a fixed block.} Fix a block, and write
\[
    s=a_{k-1}+1,\qquad t=a_k.
\]
If $s=t$, then the block is one-dimensional, and there are no root subgroups to find in it. Assume $s<t$. Since there is no $H$-invariant standard flag subspace strictly between $V_{a_{k-1}}$ and $V_{a_k}$, the $H$-invariant subspace
\[
    V_{a_{k-1}}+\langle h\mathbf{v}_s\mid h\in H\rangle
\]
is stabilized by $H$. It contains $\mathbf{v}_s$, so it is larger than $V_{a_{k-1}}$; it is contained in $V_t$, because $H\leq P$ and $P$ stabilizes $V_t$. Therefore it must equal $V_t$. Thus there is some $h\in H$ such that the coefficient of $\mathbf{v}_t$ in $h\mathbf{v}_s$ is non-zero.\\
Write $h=b_1wb_2$ by the \nameref{thm:bruhat}, with $b_1,b_2\in B$ and $w\in W$. Since $B\leq H$ and $h\in H$, we also have $w\in H$. Since $H\leq P$, the element $w$ lies in $P$, so it permutes the indices inside each block of the staircase. The row-pivot description in Corollary~\ref{cor:4.10}, applied to the $s$th column, gives
\[
    w(s)=t.
\]
Indeed, the first $s-1$ columns of $h$ lie in $V_{s-1}$, while the $s$th column lies in $V_t$ and has a non-zero $\mathbf{v}_t$-coefficient. Thus the row chosen by the Bruhat pivot construction for column $s$ is row $t$.\\
Let $q=w^{-1}(s)$. Because $w$ preserves the block $\{s,\dots,t\}$ and sends $s$ to $t$, we have
\[
    s<q\leq t.
\]
Hence $X_{sq}\leq B\leq H$. By Lemma~\ref{lem:4.6}(v),
\[
    wX_{sq}(\alpha)w^{-1}=X_{ts}(\alpha)
\]
for every $\alpha\in F$. Therefore
\[
    X_{ts}\leq H.
\]
\textit{Step 3: use commutators to get all root subgroups in the block.} The upper root subgroups $X_{ij}$ with $s\leq i<j\leq t$ already lie in $B$, and hence in $H$. Together with $X_{ts}\leq H$, the commutator formula in Lemma~\ref{lem:4.6}(iv) gives the remaining root subgroups inside the block. For $s<i<t$ and $\alpha\in F$,
\[
    X_{ti}(\alpha)=[X_{ts}(\alpha),X_{si}(1)]\in H
\]
and
\[
    X_{is}(\alpha)=[X_{it}(\alpha),X_{ts}(1)]\in H.
\]
For distinct $s<i,j<t$,
\[
    X_{ij}(\alpha)=[X_{is}(\alpha),X_{sj}(1)]\in H.
\]
Thus $H$ contains every root subgroup $X_{ij}$ supported inside this block. Since $H$ also contains all diagonal matrices, Theorem~\ref{thm:4.12} applied to this block shows that $H$ contains the full general linear group on the block.\\
\textit{Step 4: finish block by block.} Repeating Steps 2 and 3 for every block of the staircase shows that every block factor of $L_P$ is contained in $H$. Hence
\[
    L_P\leq H.
\]
Since $U_P\leq H$ as well, we have
\[
    P=U_P\rtimes L_P\leq H.
\]
We already had $H\leq P$, so $H=P$. Therefore every subgroup of $G$ containing $B$ is a staircase group, which completes the general case of Theorem~\ref{thm:5.4}.
\end{solution}

\begin{exercise}
\label{ex:5.8}
Assume that the case $r = 1$ of Theorem~\ref{thm:5.4} holds; we sketch an alternate proof of the general case of Theorem~\ref{thm:5.4}. We use induction on $n$, where $G = \mathrm{GL}(n, F)$. Let $H$ be a subgroup of $G$ which contains $B$, and suppose that $H$ stabilizes $V_m$, where $1 \leq m < n$. Let $S$ be the stabilizer of $V_m$; observe that $H \leq S$. We can write $S = U \rtimes (K_1 \times K_2)$, where the elements of $K_1 \cong \mathrm{GL}(m, F)$ have their only non-zero entries in the first $m$ rows and columns, and where $K_2$ is a direct product of general linear groups whose elements have their non-zero entries concentrated in the last $n - m$ rows and columns. Since $U \leq B \leq H$, by Exercise~\ref{ex:2.10} we have $H = U \rtimes Q$ for some $Q \leq K_1 \times K_2$. Use Goursat's theorem (Exercise~\ref{ex:2.8}), Exercise~\ref{ex:5.6}, and the induction hypothesis to analyze $Q$ and thereby conclude that $H$ is a staircase group.
\end{exercise}
\begin{solution}
We prove the general case of Theorem~\ref{thm:5.4} by strong induction on $n$, assuming that the case $r=1$ is already known. Here the case $r=1$ means the following: if $B\leq H\leq \GL(n,F)$ and $H$ stabilizes no non-zero proper standard flag subspace, then $H=\GL(n,F)$. The case $n=1$ is immediate, so assume $n>1$ and that the result is known for all smaller dimensions.\\
Let $H$ be a subgroup with $B\leq H\leq G$. If $H$ stabilizes no proper standard flag subspace $V_m$ with $1\leq m<n$, then the assumed $r=1$ case gives $H=G$, which is a staircase group. Thus we may assume that $H$ stabilizes some $V_m$, where $1\leq m<n$. Let $S$ be the stabilizer of $V_m$. Then $H\leq S$. With respect to the decomposition
\[
    V_n(F)=V_m\oplus \langle \mathbf{v}_{m+1},\dots,\mathbf{v}_n\rangle,
\]
the elements of $S$ have block form
\[
    \begin{pmatrix}
        * & *\\
        0 & *
    \end{pmatrix}.
\]
Let $U$ be the subgroup with identity diagonal blocks and arbitrary upper-right block, let $K_1\cong \GL(m,F)$ be the upper-left diagonal block group, and let $K_2\cong \GL(n-m,F)$ be the lower-right diagonal block group. Then
\[
    S=U\rtimes (K_1\times K_2).
\]
Also $U\leq B\leq H$. By Exercise~\ref{ex:2.10}, if a subgroup of a semidirect product contains the normal factor, then it is the semidirect product of that normal factor with its intersection with the complement. Applying this to $U\leq H\leq S$, we get
\[
    H=U\rtimes Q,
    \qquad
    Q=H\cap (K_1\times K_2).
\]
Thus the problem has been reduced to understanding the subgroup $Q$ of the block diagonal group $K_1\times K_2$.\\
\textit{Step 1: $Q$ contains the two block Borels.} Let
\[
    B_1=B\cap K_1,
    \qquad
    B_2=B\cap K_2.
\]
These are the standard Borel subgroups inside $K_1$ and $K_2$. Since every block diagonal element of $B_1\times B_2$ is also an element of $B$, we have
\[
    B_1\times B_2\leq B\leq H.
\]
But $B_1\times B_2\leq K_1\times K_2$, so
\[
    B_1\times B_2\leq H\cap (K_1\times K_2)=Q.
\]
\textit{Step 2: apply Goursat's theorem to $Q$.} By Exercise~\ref{ex:2.8}, Goursat's theorem says that a subgroup of a direct product is controlled by its two projections, its two coordinate kernels, and an isomorphism between the corresponding quotient groups. In our case, define
\[
    P_1=\{x\in K_1\mid (x,y)\in Q \text{ for some } y\in K_2\},
    \qquad
    I_1=\{x\in K_1\mid (x,1)\in Q\},
\]
and define $P_2$ and $I_2$ similarly. Then Exercise~\ref{ex:2.8} gives
\[
    I_i\trianglelefteq P_i\leq K_i
    \qquad
    (i=1,2),
\]
together with an isomorphism
\[
    P_1/I_1\cong P_2/I_2.
\]
Moreover, $Q$ is recovered from these data by
\[
    Q=\{(x,y)\in P_1\times P_2\mid \varphi(xI_1)=yI_2\}.
\]
The inclusion $B_1\times B_2\leq Q$ gives more than just $B_i\leq P_i$. Indeed, if $b_1\in B_1$, then $(b_1,1)\in Q$, so $b_1\in I_1$. Similarly, if $b_2\in B_2$, then $(1,b_2)\in Q$, so $b_2\in I_2$. Hence
\[
    B_i\leq I_i\leq P_i\leq K_i
    \qquad
    (i=1,2).
\]
\textit{Step 3: use induction and self-normalizing parabolics.} Since $K_1\cong \GL(m,F)$ and $K_2\cong \GL(n-m,F)$ have smaller dimension than $G$, the induction hypothesis applies inside both blocks. Because $P_i$ and $I_i$ contain the standard Borel subgroup $B_i$, the induction hypothesis tells us that $P_i$ and $I_i$ are staircase groups inside $K_i$. In particular, $I_i$ is a parabolic subgroup of $K_i$.\\
Now $I_i\trianglelefteq P_i$. Therefore every element of $P_i$ normalizes $I_i$, so
\[
    P_i\leq N_{K_i}(I_i).
\]
By Exercise~\ref{ex:5.6}, parabolic subgroups are self-normalizing. Applying this inside $K_i$, we get
\[
    N_{K_i}(I_i)=I_i.
\]
Therefore
\[
    P_i\leq I_i.
\]
Since $I_i\leq P_i$ already, we conclude that
\[
    P_i=I_i
    \qquad
    (i=1,2).
\]
\textit{Step 4: the Goursat coupling disappears.} Since $P_i=I_i$, both quotient groups in Goursat's theorem are trivial:
\[
    P_1/I_1=1,
    \qquad
    P_2/I_2=1.
\]
Thus the condition $\varphi(xI_1)=yI_2$ in the description of $Q$ is automatic. Hence
\[
    Q=P_1\times P_2.
\]
The induction hypothesis also tells us that $P_1$ and $P_2$ are staircase groups inside their respective blocks.\\
\textit{Step 5: combine the two block staircases.} Suppose $P_1$ stabilizes a standard subflag inside $V_m$, say
\[
    0\subset V_{a_1}\subset \dots\subset V_{a_r}=V_m.
\]
For the lower block, write
\[
    W_j=\langle \mathbf{v}_{m+1},\dots,\mathbf{v}_{m+j}\rangle
    \qquad
    (1\leq j\leq n-m).
\]
Thus $W_1\subset \dots\subset W_{n-m}$ is the standard flag inside the lower block. Suppose $P_2$ stabilizes the lower-block subflag
\[
    0\subset W_{b_1}\subset \dots\subset W_{b_s}=W_{n-m}.
\]
When these lower-block spaces are placed back inside $V_n(F)$, we must add the first block $V_m$. Thus
\[
    V_m\oplus W_{b_j}=V_{m+b_j}
    \qquad
    (1\leq j\leq s).
\]
Thus the two block flags combine to the standard subflag
\[
    0\subset V_{a_1}\subset \dots\subset V_m
      \subset V_{m+b_1}\subset \dots\subset V_n(F).
\]
The stabilizer of this combined subflag is exactly the group of block upper triangular matrices whose upper-left diagonal block lies in $P_1$ and whose lower-right diagonal block lies in $P_2$:
\[
    U\rtimes(P_1\times P_2).
\]
But this group is $H$, since
\[
    H=U\rtimes Q=U\rtimes(P_1\times P_2).
\]
Therefore $H$ is a staircase group. This completes the induction step, and hence proves the general case of Theorem~\ref{thm:5.4} from the assumed $r=1$ case.
\end{solution}

\end{document}
